\chapter{绪\quad 论}

\section{研究背景与意义}

平台化传播正在重组公众人物的职业生命周期。过去，艺人的可见性主要由杂志封面、电视节目、影视作品和广告投放共同塑造；现在，短视频、长视频、社交平台和自营品牌又把表达频率、互动密度与叙事控制权推到更重要的位置。职业转型因此不再只是“从模特到演员”或“从电视到网络”的岗位变化，而是平台规则、角色边界、公众期待和危机回应的连续调整。

佐佐木希（官方日文名为佐々木希，英文名 Nozomi Sasaki）是一个适合公开模板演示的案例。TopCoat 官方资料显示，她出生于 1988 年 2 月 8 日，出身地为日本秋田县\cite{topcoat-profile}。PR TIMES 2024 年发布的官方频道开设新闻进一步说明，她 2006 年进入演艺行业，此后在电影、电视剧、广告、杂志等领域持续活动，并在 2024 年开设 YouTube 官方频道“佐々木 希(仮)”\cite{youtube-prtimes}。这些事实提供了清晰的跨媒介路径，也避免了示例论文依赖真实访谈、内部数据或保密材料。

本文选择佐佐木希，不是为了复述人物八卦，而是为了展示一个长期公开职业样本如何被转化为博士论文中的研究对象：哪些事实可以进入正文，哪些敏感事件只能作为舆论背景，哪些图表可以帮助用户理解模板的排版能力。对于 \texttt{thesis-nwu-doctor} 项目而言，这种公开案例比空白示例更接近真实写作，也更便于持续回归验证。

\section{研究现状与问题意识}

职业发展理论强调个体会在探索、建立、维持与再适应之间移动\cite{super-career}。在文娱产业中，这种移动常常与媒介环境变化同步发生：纸媒时代的集中曝光强调形象识别，影视与综艺阶段强调角色扩展，社交平台阶段则强调持续表达和关系维护。融合文化研究也指出，内容在不同媒介之间流动时，生产者与受众之间的关系会被重新组织\cite{jenkins-convergence}。

形象管理研究提供了另一个切口。Goffman 关于自我呈现的讨论提醒我们，公众形象并不是固定标签，而是在具体情境中被持续表演和解释的结果\cite{goffman-presentation}。当社交平台成为职业系统的一部分时，公众人物需要更频繁地解释自我，也更容易因突发事件进入被动叙事。2020 年，围绕佐佐木希配偶渡部建的婚外情报道引发舆论事件，佐佐木希随后公开回应并表示将与配偶沟通、继续工作\cite{nikkan-2020}。本文将该事件视为职业叙事中的外部冲击，而不展开私人生活判断。

\section{研究问题}

围绕上述背景，本文提出三个问题：

\begin{enumerate}
  \item 佐佐木希的公开职业路径可以划分为哪些相对稳定的跨媒介阶段？
  \item 平面媒体、影视商业、社交平台和 YouTube 等渠道分别提供了什么样的叙事控制权？
  \item 在家庭相关舆论冲击后，公众人物如何通过内容节奏、平台选择和角色边界调整维持形象韧性？
\end{enumerate}

\section{研究设计与章节安排}

本文采用“公开事实核查 + 阶段化案例分析 + 模板能力展示”的写法。事实层面，只纳入官方网页、公开新闻和可复核的理论文献；分析层面，将案例拆解为阶段、平台、角色和回应四类变量；排版层面，则通过图、表、公式、参考文献、附录、成果和作者简介完整展示西北大学博士论文模板的常见结构。

全文安排如下：第一章说明研究背景、事实边界和研究问题；第二章建立职业阶段与形象韧性分析框架；第三章结合复用图片、表格和公式展示案例分析结果与模板要素。后置部分依次给出结论、参考文献、附录、攻读博士学位期间取得的研究成果、致谢和作者简介。
